\chapter{\MakeUppercase{Introdução}}

Neste capítulo introdutório, serão apresentados o contexto do projeto desenvolvido durante o estágio supervisionado, a empresa cliente e o cenário do agronegócio que motivou a construção da plataforma de protocolos. Serão também delineados os objetivos do estágio e deste relatório, que visa documentar as atividades realizadas, as tecnologias adotadas e os resultados alcançados.

\section{Contextualização}
O agronegócio é um dos setores mais importantes e dinâmicos da economia, exigindo constante inovação para maximizar a produtividade e a sustentabilidade no campo. Dentro desse cenário, a pesquisa agrícola focada na análise de solos e sementes desempenha um papel fundamental. Com o avanço da chamada Agricultura 4.0, a coleta, o armazenamento e a análise de dados tornaram-se indispensáveis para a tomada de decisões precisas.

Apesar da crescente digitalização, muitas instituições de pesquisa ainda enfrentam o desafio da fragmentação de dados, dependendo de ferramentas manuais e planilhas eletrônicas para registrar suas descobertas. Esse cenário dificulta a rastreabilidade das informações, a aplicação de regras de negócio complexas e, principalmente, a geração de análises visuais rápidas que poderiam acelerar os resultados das pesquisas. É nesse contexto de modernização e centralização de dados que se insere o projeto desenvolvido durante o estágio.

\section{A Empresa Cliente e o Cenário do Projeto}
As atividades de estágio supervisionado foram realizadas na empresa Zeeway, uma empresa lavrense de tecnologia fundada no ecossistema de inovação da Universidade Federal de Lavras (UFLA). A Zeeway atua no desenvolvimento de software e na entrega de soluções tecnológicas personalizadas, destacando-se pelo compromisso com a qualidade dos projetos e com os resultados dos seus clientes. Atualmente, a empresa conta com um quadro de 15 a 30 colaboradores diretos e mantém uma base de aproximadamente 40 clientes diretos e 500 clientes indiretos. O trabalho principal desenvolvido durante o período de estágio foi a construção de uma plataforma web encomendada pela Terra Gerais, uma empresa do setor do agronegócio especializada em pesquisa agrícola de solos e sementes.

A necessidade da Terra Gerais consistia em abandonar o uso descentralizado de planilhas e adotar uma plataforma própria para registrar os seus "protocolos" — registros detalhados com todas as informações de uma safra, componentes do solo e produtividade de grãos. O acesso ao sistema seria distribuído aos colaboradores, exigindo uma arquitetura robusta para lidar com diferentes perfis de usuário e permissões de acesso.

A participação no projeto concentrou-se no desenvolvimento da interface visual e da experiência do usuário (front-end). A aplicação foi construída utilizando a biblioteca React com TypeScript, empregando componentes da interface Material UI (MUI) e integrando a biblioteca ApexCharts para o grande diferencial do sistema: a geração automatizada de gráficos de análise agronômica. Todo o ciclo de desenvolvimento foi pautado pela metodologia ágil Scrum, permitindo organização e entregas incrementais.

\section{Objetivos do Estágio e do TCC}

O objetivo principal do estágio foi aplicar, em um ambiente corporativo real, os conhecimentos adquiridos ao longo da graduação, vivenciando a rotina de desenvolvimento de software em equipe e os desafios da engenharia de requisitos.

Este relatório tem como objetivo apresentar de forma estruturada as atividades desempenhadas na Zeeway, detalhando a arquitetura front-end construída, as tecnologias adotadas e a metodologia de trabalho. Busca-se demonstrar também os desafios técnicos enfrentados na tradução de regras de negócio complexas para a interface da plataforma e o impacto da experiência na formação profissional e acadêmica.